\subsection{Clasificación}

Según el tipo de organización y las leyes vigentes, la Universidad Nacional de Ingeniería se clasifica de la siguiente manera:

\begin{itemize}
    \item \textbf{Según su propiedad y régimen legal:} Es una institución \textbf{pública} que cuenta con autonomía académica, administrativa, económica y de gobierno de acuerdo con la Ley Universitaria y la Constitución Política del Perú \cite{ley30220}.
    
    \item \textbf{Según su sector y actividad económica:} Pertenece al sector terciario (servicios), ya que brinda \textbf{educación superior universitaria e investigación}. Sus actividades principales son la enseñanza, la investigación científica y el desarrollo tecnológico.
    
    \item \textbf{Según su finalidad y régimen de utilidades:} Es una institución \textbf{sin fines de lucro}. Todos los recursos presupuestales provenientes del tesoro público y de sus recursos directamente recaudados se reinvierten en el cumplimiento de sus objetivos institucionales, infraestructura y bienestar de la comunidad universitaria.
    
    \item \textbf{Según su tamaño y alcance:} Es una \textbf{gran organización} de alcance \textbf{nacional}. Cuenta con un campus central en Lima que reúne a once facultades y recibe a postulantes y estudiantes de todas las regiones del país.
\end{itemize}
